\section{Scale Sensitivity in Self-Gated Activations}\label{sec:problem}

%% ============================================================
\subsection{The self-gated family}\label{sec:selfgated}
%% ============================================================

A self-gated activation takes the form $f(z_i) = z_i \cdot g(z_i)$, where $g: \mathbb{R} \to (0,1)$ is a smooth, monotonic gate.
The neuron gates itself: the same scalar $z_i$ determines both the signal and the gating decision.
The dominant instances are GELU ($g = \Phi$, the standard-normal CDF) and SiLU ($g = \sigma$, the logistic sigmoid), both widely adopted in modern deep learning~\citep{devlin2019bert,dosovitskiy2021image,brown2020language}.

Compared to ReLU, these activations provide smooth gradients, graded modulation (gate values between 0 and 1, rather than binary), and no dead neurons.
All three advantages stem from replacing the binary indicator $\mathbf{1}[z_i > 0]$ with a smooth function $g$.
We now show that this replacement has a side effect.

%% ============================================================
\subsection{The scale invariance trade-off}\label{sec:tradeoff}
%% ============================================================

An activation $f$ is \emph{positively homogeneous} if $f(\alpha\mathbf{z}) = \alpha\,f(\mathbf{z})$ for all $\alpha > 0$, which is equivalent to the output direction $f(\mathbf{z})/\|f(\mathbf{z})\|$ being independent of scale.

Piecewise-linear activations---ReLU, Leaky ReLU, PReLU---are positively homogeneous, since $\max(\alpha z, 0) = \alpha\max(z, 0)$.
In the self-gated form, however, positive homogeneity requires $g(\alpha z_i) = g(z_i)$ for all $\alpha$:

\begin{proposition}\label{prop:no_invariance}
No self-gated activation $z_i \cdot g(z_i)$ with non-constant smooth $g$ is positively homogeneous.
\end{proposition}
\begin{proof}
Positive homogeneity requires $\alpha z_i \cdot g(\alpha z_i) = \alpha \cdot z_i \cdot g(z_i)$, hence $g(\alpha z_i) = g(z_i)$ for all $\alpha > 0$ and $z_i \neq 0$.
This forces $g$ to be constant on $(0,\infty)$ and on $(-\infty, 0)$; continuity then makes $g$ constant everywhere, contradicting the definition of a gate.
\end{proof}

Smoothness and scale invariance are thus incompatible within the self-gated form.
Every smooth self-gated activation---GELU, SiLU, and others---necessarily lost the scale invariance that ReLU had.

%% ============================================================
\subsection{Output direction entanglement}\label{sec:direction}
%% ============================================================

What is the concrete consequence of losing scale invariance?
We decompose any pre-activation vector as $\mathbf{z} = \rho\,\hat{\mathbf{z}}$, where $\rho = \mathrm{rms}(\mathbf{z})$ and $\hat{\mathbf{z}} = \mathbf{z}/\rho$.

\begin{proposition}\label{prop:direction}
\textbf{(Output direction.)}

\emph{(i)} For any positively homogeneous activation: $\hat{\mathbf{y}}(\alpha\mathbf{z}) = \hat{\mathbf{y}}(\mathbf{z})$ for all $\alpha > 0$.

\emph{(ii)} For a self-gated activation $f(z_i) = z_i \cdot g(z_i)$ with non-constant $g$:
\begin{equation}
    \hat{\mathbf{y}}(\mathbf{z}) = \frac{\hat{\mathbf{z}} \odot g(\rho\hat{\mathbf{z}})}{\|\hat{\mathbf{z}} \odot g(\rho\hat{\mathbf{z}})\|},
\end{equation}
which depends on $\rho$.
Two inputs sharing the same pattern $\hat{\mathbf{z}}$ but differing in scale $\rho$ produce different output directions.

\emph{(iii)} For GELU, as $\rho \to \infty$:
$\hat{\mathbf{y}}(\mathbf{z}) \to \mathrm{ReLU}(\hat{\mathbf{z}})/\|\mathrm{ReLU}(\hat{\mathbf{z}})\|$.
The smooth gate saturates to a binary gate, and the output direction degenerates to ReLU's.
\end{proposition}

\begin{proof}
(i) $f(\alpha\mathbf{z}) = \alpha f(\mathbf{z})$; dividing by the norm cancels $\alpha$.
(ii) $f(\mathbf{z})_i = \rho\hat{z}_i \cdot g(\rho\hat{z}_i)$; the direction is $\hat{z}_i g(\rho\hat{z}_i)$ normalized, which depends on $\rho$ when $g$ is non-constant.
(iii) $\Phi(\rho\hat{z}_i) \to \mathbf{1}[\hat{z}_i > 0]$ as $\rho \to \infty$.
\end{proof}

\paragraph{Why direction matters.}
As noted in Section~\ref{sec:intro}, the next layer decomposes its response as $\langle \mathbf{w}_j, \mathbf{y} \rangle = \|\mathbf{y}\| \cdot \langle \mathbf{w}_j, \hat{\mathbf{y}} \rangle$: scale acts as gain, direction determines selectivity.
In the self-gated form, a common scale factor changes both, entangling the two roles.

%% ============================================================
\subsection{Consequences for learning}\label{sec:consequences}
%% ============================================================

\paragraph{A harder optimization problem.}
With a scale-invariant activation, the next layer's weight $\mathbf{w}_j$ must learn a function of $\hat{\mathbf{z}}$ alone.
With a scale-dependent one, it must learn a function of $(\hat{\mathbf{z}}, \rho)$ jointly---a strictly larger input space.
If the test-time distribution of $\rho$ differs from training, the learned weights encounter output directions they were not optimized for.

\paragraph{Training dynamics compound the issue.}
The output of a self-gated activation increases with scale:
\begin{equation}\label{eq:df_dalpha}
    \frac{\partial}{\partial \alpha}\bigl[\alpha \tilde{z}_i\,g(\alpha \tilde{z}_i)\bigr]
    = \tilde{z}_i\,g(\alpha \tilde{z}_i) + \alpha\,\tilde{z}_i^2\,g'(\alpha \tilde{z}_i) \geq 0 \quad (\tilde{z}_i > 0),
\end{equation}
so cross-entropy (which rewards larger logits) can drive weight norms upward.
In a piecewise-linear activation, this only amplifies the output; the direction is unchanged.
In a smooth self-gated activation, increasing $\alpha$ simultaneously shifts the output direction (Proposition~\ref{prop:direction}(ii)) and sharpens the gate toward saturation (Proposition~\ref{prop:direction}(iii)).
A single optimization pressure thus has a compound side effect: amplification \emph{and} direction drift.

\paragraph{LayerNorm does not prevent this.}
In a Transformer MLP, $\mathbf{z} = W\,\mathrm{LN}(\mathbf{x}) + \mathbf{b}$.
LayerNorm controls $\|\mathbf{x}\|$, but $W$ reintroduces scale: $\mathrm{Var}(z_i) \propto \|\mathbf{w}_i\|^2$.
The scale sensitivity arises \emph{inside the activation}, after the projection; inter-layer normalization cannot reach it (Appendix~\ref{app:ln_ablation}).